\section{Problem Statement}

Design and implement a memory allocator that simulates a heap with explicit free list management and boundary tag coalescing. The allocator must support \texttt{malloc(size)} and \texttt{free(ptr)} operations and implement first-fit, best-fit, and next-fit allocation strategies. The system must track fragmentation metrics and correctly handle splitting of large free blocks and coalescing of adjacent free blocks upon deallocation.

\section{Core Concepts}

\begin{itemize}
  \item \textbf{Block Header:} Metadata stored at the beginning of each block, recording the block size and whether the block is currently allocated or free.
  \item \textbf{Block Footer:} A copy of the header placed at the end of the block. The footer enables backward traversal, allowing the allocator to inspect the preceding block during coalescing.
  \item \textbf{Free List:} A linked list threaded through the free blocks in the heap. Each free block contains pointers to the next and previous free blocks, enabling efficient traversal without scanning the entire heap.
  \item \textbf{Splitting:} When a free block is larger than the requested allocation, it is divided into an allocated block of the requested size and a smaller remaining free block.
  \item \textbf{Coalescing:} When a block is freed, the allocator checks the adjacent blocks. If a neighbor is also free, the two blocks are merged into a single larger free block to reduce fragmentation.
\end{itemize}

\section{Key Challenges}

\begin{itemize}
  \item \textbf{Pointer Arithmetic:} Navigating between headers, payloads, and footers requires precise offset calculations. Off-by-one errors in pointer arithmetic lead to heap corruption.
  \item \textbf{Boundary Tags:} Using footers to enable backward coalescing introduces the invariant that every block must maintain a consistent header-footer pair, even after splitting.
  \item \textbf{Free List Maintenance:} Updating the doubly-linked free list pointers when blocks are split, coalesced, allocated, or freed, ensuring no dangling or circular references.
  \item \textbf{Minimum Block Size:} Free blocks must be large enough to store the list pointers in addition to the header and footer, imposing a minimum allocation granularity.
  \item \textbf{Fragmentation Measurement:} Computing external fragmentation as the ratio of unusable scattered free memory to total free memory, and reporting it after each operation.
\end{itemize}

\section{Sample Input}

\begin{lstlisting}[style=json, caption={data/input\_sample.json}]
{
  "heap_size": 1024,
  "allocation_strategy": "first_fit",
  "operations": [
    {"type": "malloc", "id": "A", "size": 128},
    {"type": "malloc", "id": "B", "size": 256},
    {"type": "malloc", "id": "C", "size": 64},
    {"type": "free", "id": "A"},
    {"type": "free", "id": "C"},
    {"type": "malloc", "id": "D", "size": 100},
    {"type": "heap_dump"},
    {"type": "free", "id": "B"},
    {"type": "heap_dump"},
    {"type": "malloc", "id": "E", "size": 400}
  ]
}
\end{lstlisting}

\vfill
\noindent\rule{\textwidth}{0.4pt}
{\small\textit{For full project requirements, STL restrictions, submission format, and grading criteria, refer to \textbf{base.pdf} (available on the \href{https://github.com/BestSithInEU/cse211-term-projects/releases}{Releases page}).}}
